\chapter{Implementation}

This chapter describes the Java implementation of the Espresso algorithm. It covers the data structures used, the BLIF file parsing framework, the build system, and the implementation of the core algorithmic steps: Complement, Expand, Irredundant, and Reduce. Java was chosen for its speed and Object Oriented abstractions.

\section{Data Structure}
\label{sec:datastructure}

The first thing to decide is the choice of data structure for representing Boolean covers. \ac{PCN} is used as the cube representation throughout the implementation. Each literal within a cube is encoded as a 2-bit integer: \texttt{1} (\texttt{0b01}) represents TRUE, \texttt{2} (\texttt{0b10}) represents FALSE, and \texttt{3} (\texttt{0b11}) represents a \ac{DC} condition. This encoding maps directly to efficient bitwise AND and OR operations, which are used extensively throughout the Espresso sub-parts.

The cover is stored as \texttt{List$<$int[]$>$}, backed by a \texttt{LinkedList}, where each \texttt{int[]} array represents a single cube. The index \texttt{cover[i][j]} holds the PCN state of the $j$-th literal in the $i$-th cube. The number of columns (inputs) is fixed for a given function instance, while the number of rows (cubes) varies as the algorithm inserts and removes cubes during minimization. A \texttt{LinkedList} is used because insertions and deletions are faster compared to fixed arrays, which is useful for frequent cube manipulation required by Expand, Irredundant, and Reduce.

As an example, consider the function $F = abc + a'b + c'$, represented in PCN as:

\[
\begin{pmatrix}
1 & 1 & 1 \\
2 & 1 & 3 \\
3 & 3 & 2
\end{pmatrix}
\]

where rows correspond to cubes and columns correspond to literals $a$, $b$, and $c$ respectively.

Additional supporting data structures include primitive \texttt{int[]} arrays for storing literal ordering and \texttt{ArrayList} for sorting literals by weight heuristics during binate variable selection and, ordering during expansion and reduction.

\section{Parsing Framework}

The input to the tool is a \texttt{.blif} file in \ac{BLIF}. Parsing is handled by the \texttt{Blif} class, which reads the file line by line using a \texttt{BufferedReader}. Each \texttt{.names} directive marks the start of a new function block; the subsequent cube lines are parsed into \ac{PCN} \texttt{Cube} objects via a character-by-character switch: \texttt{0} maps to FALSE (\texttt{0b10}), \texttt{1} maps to TRUE (\texttt{0b01}), and \texttt{-} maps to \ac{DC} (\texttt{0b11}). Only on-set cubes (output column \texttt{1}) are supported. The parser also preserves the file preamble so that the minimized result can be written back to a new \texttt{.blif} file in an \texttt{output/} subdirectory, with the cover replaced by the minimized one.

\section{Build System}

The project is built with \texttt{make}. A suffix rule \texttt{.java.class} compiles each source file via \texttt{javac -g}, and the \texttt{classes} target drives compilation of all nine source files in dependency order. Running \texttt{make} from the \texttt{src/} directory produces the \texttt{.class} files; \texttt{make clean} removes them. The tool is then invoked as \texttt{java Main <file.blif>}.

\section{Espresso Implementation}

The \texttt{Espresso} class is the main class of the minimization algorithm. It is instantiated with the list of \texttt{Function} objects produced by the parser. The constructor first converts the bitset-based \texttt{Cube} representation into the \texttt{List$<$int[]$>$} cover format described in Section~\ref{sec:datastructure}, then instantiates one object each of \texttt{Complement}, \texttt{Expand}, \texttt{Irredundant}, and \texttt{Reduce}, passing \texttt{numInputs} to each.

The algorithm has two phases. In the pre-processing phase, the complement of the initial cover is computed once via recursive Shannon expansion (\texttt{Complement}). The cover is then expanded (\texttt{Expand}), which raises literals to \ac{DC} as far as possible without intersecting the complement, followed by a removal of redundant cubes (\texttt{Irredundant}).

The cost function used throughout is the \textbf{number of cubes} $|F|$; minimization terminates when no iteration reduces this count. The main loop iterates Reduce $\to$ Expand $\to$ Irredundant until the cover size no longer decreases:

\begin{algorithm}
\caption{Espresso main loop}
\begin{algorithmic}[1]
\Repeat
    \State $\text{prevSize} \gets |\text{cover}|$
    \State $\text{cover} \gets \textsc{Reduce}(\text{cover})$
    \State $\text{cover} \gets \textsc{Expand}(\text{cover},\, \text{complement})$
    \State $\text{cover} \gets \textsc{Irredundant}(\text{cover})$
\Until{$|\text{cover}| \geq \text{prevSize}$}
\end{algorithmic}
\end{algorithm}

\noindent The minimized cover is stored as \texttt{minimizedCover} and returned via \texttt{getMinimizedCover()}, after which \texttt{Main} converts it back to \texttt{Cube} objects and writes the result to a \texttt{.blif} file.

\section{Complement}

The \texttt{Complement} class computes $F'$ via recursive Shannon expansion. The result is passed to \texttt{Expand} as the off-set, ensuring no conflict between ON-set and Off-set. The three base cases are: an empty cover returns the universe; a cover containing the universal cube returns empty; a single cube applies De Morgan's theorem directly (each TRUE/FALSE literal becomes its own all-\ac{DC} output cube with that literal flipped). Otherwise, a split variable $x$ is chosen, the positive and negative cofactors are complemented recursively, and the results are merged as $F' = x' \cdot F_{x=1}' + x \cdot F_{x=0}'$.

The splitting order is based on  \textbf{variable ordering heuristic}. Before recursing, \texttt{getOrder} counts how often each variable appears as TRUE vs.\ FALSE across all cubes. Variables that appear in both states (\emph{binate}) are split on first, since they reduce the cover size in both cofactors simultaneously. Within that group, variables are sorted by $\min(\text{trueCount}, \text{falseCount})$ descending, favouring the most balanced split. Unate variables follow. This keeps the recursion tree balanced and avoids deep sub-tree depth.

\section{Expand}

The \texttt{Expand} class raises each literal in each cube to \ac{DC} as far as possible without entering the off-set. For each literal, it is tentatively raised to \ac{DC} and \texttt{checkIntersect} verifies the result against every cube in the complement cover: two cubes intersect if and only if their PCN values have a non-zero bitwise AND in every column. If any complement cube intersects, the literal is restored.

The \textbf{weight heuristic} in \texttt{getOrder} orders cubes by ascending dot product of a row vector (the cube's PCN bits) and a column vector (the per-literal TRUE/FALSE occurrence counts across the whole cover). Cubes with a low dot product share few literals with other cubes and are therefore expanded first, they are less likely to block each other's expansion.

\section{Irredundant}

The \texttt{Irredundant} class drops cubes that contribute no unique minterms. For each cube $c$, the cover $F \setminus \{c\}$ is formed and its complement $(F \setminus \{c\})'$ is computed. Cube $c$ is redundant if and only if $c \cap (F \setminus \{c\})' = \emptyset$ - meaning every minterm of $c$ is already covered by the rest. The intersection check reuses the same bitwise test as Expand. Redundant cubes are removed in place.

\section{Reduce}

The \texttt{Reduce} class shrinks each cube to its smallest cover that still holds at least one minterm unique to it, preparing for a different Expand pass. For cube $c_i$, the \emph{supercube} is computed as the bitwise OR of all cubes in $(F \setminus \{c_i\})'$, i.e.\ the region not already covered by the other cubes. The reduced cube is then $c_i \cap \text{supercube}$, computed with a bitwise AND. This guarantees $c_i$ captures minterms that no other cube covers.

The weight heuristic mirrors Expand but sorts in \emph{descending} order, so the most-shared cubes (highest dot product) are reduced first since they have more overlap with others and can be shrunk the most aggressively.\ainote{Claude Sonnet 4.6}{Improve LaTeX formatting and grammar of the Implementation chapter.}

